\documentclass[12pt,a4paper]{article}

\usepackage[margin=1in]{geometry}
\usepackage{amsmath}
\usepackage{amssymb}
\usepackage{booktabs}
\usepackage{array}
\usepackage{enumitem}
\usepackage{hyperref}
\usepackage{xcolor}

\hypersetup{
    colorlinks=true,
    linkcolor=blue,
    urlcolor=blue,
    citecolor=blue
}

\title{\textbf{Industry Measures versus Rigorous Mathematical Evaluation}\\
\large Weekly Mentor--Fellow Scheduling System}
\author{Marcelino}
\date{\today}

\begin{document}

\maketitle

\begin{abstract}
This report discusses the difference between practical industry measures and rigorous mathematical measures in the context of a weekly mentor--fellow scheduling system. The project compares multiple scheduling backends, including exhaustive baseline search, an improved heuristic-oriented backend, and an AI-generated scarcity-aware greedy backend. The main point is that industry evaluation often emphasizes usability, speed, maintainability, and operational reliability, while mathematical evaluation emphasizes formal feasibility, optimality, guarantees, and provable bounds.
\end{abstract}

\tableofcontents
\newpage

\section{Project Context}

The system schedules weekly one-on-one mentorship sessions between fellows and mentors. Each fellow request includes:

\begin{itemize}[noitemsep]
    \item a fellow identifier,
    \item a requested topic,
    \item an urgency level,
    \item a set of available slot identifiers,
    \item a week or timestamp value.
\end{itemize}

Each mentor has a set of specialty topics, and each slot belongs to exactly one mentor. A valid assignment must satisfy topic compatibility, fellow availability, slot capacity, and the one-slot-per-fellow rule.

\section{Mathematical View of the Problem}

\subsection{Formal Inputs}

Let:

\begin{itemize}[noitemsep]
    \item $F$ be the set of fellows,
    \item $M$ be the set of mentors,
    \item $S$ be the set of available slots,
    \item $R$ be the set of weekly requests.
\end{itemize}

Each request $r \in R$ has a topic, urgency weight, available slot set, and time value. Each mentor $m \in M$ has a set of specialties.

\subsection{Feasibility Constraints}

An assignment is feasible only if:

\begin{align}
    \sum_{s \in S} x_{r,s} &\leq 1 && \forall r \in R, \\
    \sum_{r \in R} x_{r,s} &\leq 1 && \forall s \in S, \\
    x_{r,s} &= 0 && \text{if } s \notin A_r, \\
    x_{r,s} &= 0 && \text{if mentor}(s) \text{ does not cover topic}(r).
\end{align}

\subsection{Objective Function}

The mathematical objective can be expressed as a minimization problem over unserved or delayed requests. One common form is:

\begin{equation}
    \text{Penalty}_{week}
    =
    \sum_{f \in U}
    w(\text{urgency}_f)(s_f + 1)^2,
\end{equation}

where $U$ is the set of unserved fellows, $w$ is the urgency weight, and $s_f$ is the starvation counter.

\subsection{Rigorous Mathematical Questions}

A mathematical analysis asks questions such as:

\begin{itemize}[noitemsep]
    \item Is the produced schedule feasible?
    \item Is the objective function minimized?
    \item Is the approximation ratio bounded?
    \item Does the algorithm always terminate?
    \item What is the worst-case time complexity?
    \item Can a counterexample show that a heuristic is not optimal?
\end{itemize}

\section{Industry View of the Problem}

\subsection{Practical Success Measures}

Industry evaluation usually asks whether the system is useful, reliable, and fast enough for real users. Common measures include:

\begin{itemize}[noitemsep]
    \item response time,
    \item throughput,
    \item percentage of requests assigned,
    \item user-facing clarity,
    \item operational stability,
    \item ease of debugging,
    \item maintainability,
    \item deployment simplicity.
\end{itemize}

\subsection{Frontend and Product Measures}

For this project, the Python client introduces practical measures that are not purely mathematical:

\begin{itemize}[noitemsep]
    \item ability to switch algorithms during runtime,
    \item ability to reset to a known mid-lifecycle state,
    \item ability to load canned edge-case CSVs,
    \item live sanity checks for impossible requests,
    \item readable weekly schedule views,
    \item simple operational workflow for a user.
\end{itemize}

\subsection{Performance Measures}

The implemented frontend records runtime metrics such as:

\begin{itemize}[noitemsep]
    \item backend selected,
    \item elapsed wall-clock time,
    \item number of requests submitted,
    \item number of requests assigned,
    \item weekly penalty,
    \item timeout or fallback behavior.
\end{itemize}

\section{Comparison of Measurement Styles}

\begin{table}[h]
\centering
\renewcommand{\arraystretch}{1.25}
\begin{tabular}{p{0.28\linewidth} p{0.31\linewidth} p{0.31\linewidth}}
\toprule
\textbf{Aspect} & \textbf{Rigorous Mathematics} & \textbf{Industry Practice} \\
\midrule
Correctness & Formal feasibility and proof obligations & No visible failures and valid user-facing output \\
Optimality & Exact minimum or bounded approximation & Good enough assignment quality under time limits \\
Performance & Asymptotic complexity and worst-case bounds & Measured latency on real or representative cases \\
Failure Handling & Classified infeasibility and counterexamples & Clear warnings, timeouts, fallback behavior \\
Evaluation Data & Constructed proofs and adversarial cases & Regression cases, CSV scenarios, user workflows \\
System State & Mathematical variables such as $s_f$ & Reset buttons, snapshots, session state, diagnostics \\
\bottomrule
\end{tabular}
\caption{Difference between mathematical and industry-oriented evaluation.}
\end{table}

\section{Algorithm Backends in This Project}

\subsection{Baseline Backend}

The baseline backend represents the mathematically strict reference style. It searches exhaustively on smaller inputs and is useful for validating correctness and comparing heuristic behavior.

\subsection{Improved Backend}

The improved backend balances rigor and practicality. It preserves the same input contract but introduces preprocessing, timeout handling, and a fallback heuristic.

\subsection{AI Backend}

The AI backend represents a scarcity-aware greedy approach. It is useful for practical comparison because it often runs quickly and exposes a different scheduling philosophy.

\section{Why Industry Measures Differ from Mathematical Measures}

\subsection{Time Limits Matter}

In mathematics, an exponential algorithm can still be valid if it is correct. In industry, an algorithm that takes too long is usually considered unusable even if it is exact.

\subsection{Users Need Explanations}

Formal feasibility is not enough for a real interface. A user needs to know why a request cannot be scheduled, whether a topic is unknown, and whether a listed slot is incompatible.

\subsection{Systems Need Recovery Paths}

Industry systems need reset buttons, snapshots, fallback algorithms, and diagnostics. These are not part of the pure optimization model, but they are essential for testing and operation.

\subsection{Metrics Are Often Empirical}

Mathematical analysis focuses on proofs. Industry analysis often relies on measured runtime, observed behavior, logs, dashboards, and repeated test scenarios.

\section{Test Cases as a Bridge Between Both Views}

The generated CSV cases help connect mathematical reasoning with practical testing:

\begin{itemize}[noitemsep]
    \item \texttt{all\_same\_slot.csv}: capacity bottleneck,
    \item \texttt{impossible\_topic.csv}: infeasible topic and incompatible slots,
    \item \texttt{no\_overlap.csv}: trivial perfect matching,
    \item \texttt{topic\_bottleneck.csv}: low topic coverage,
    \item \texttt{starvation\_rescue.csv}: starvation versus urgency,
    \item \texttt{large\_mixed.csv}: practical timeout and performance comparison.
\end{itemize}

\section{Discussion}

The rigorous mathematical view gives the project a precise definition of correctness. It defines what a legal assignment is, how penalty is computed, and what it means for an algorithm to be optimal or near-optimal.

The industry view gives the project a usable system. It asks whether users can submit requests, understand failures, compare algorithms, reload known states, and observe runtime behavior.

Both views are necessary. Mathematical rigor prevents invalid schedules and vague objectives. Industry measures ensure the system can actually be used, tested, and demonstrated under realistic constraints.

\section{Conclusion}

This project shows that mathematical correctness and industry usefulness are related but not identical. A rigorous algorithm can be too slow for practical use, while a practical heuristic can be fast and user-friendly without offering full optimality guarantees. The best engineering result is usually a system that keeps the mathematical model explicit while exposing practical controls, diagnostics, and performance measurements.

\end{document}
